



\section{BERT4Rec}
Before going into the details, we first introduce the research problem, the basic concepts, and the notations in this paper.

\subsection{Problem Statement}
In sequential recommendation, let $\mathcal{U}{=}\{u_1,u_2,\dots,u_{\vert\mathcal{U}\vert}\}$ denote a set of users, $\mathcal{V}{=}\{v_1,v_2,\dots,v_{|\mathcal{V}|}\}$ be a set of items, and list $\mathcal{S}_u{=}[v_1^{(u)},\dots,v_t^{(u)},\dots,v_{n_u}^{(u)}]$ denote the interaction sequence in chronological order for user $u \in \mathcal{U}$, where $v_t^{(u)} \in \mathcal{V}$ is the item that $u$ has interacted with at time step\footnote{Here, following~\cite{Rendle:WWW2010:FPM,Kang:ICDM2018:SAS}, we use the relative time index instead of absolute time index for numbering interaction records.} $t$ and $n_u$ is the the length of interaction sequence for user $u$.
Given the interaction history $\mathcal{S}_u$, sequential recommendation aims to predict the item that user $u$ will interact with at time step $n_{u} + 1$.
It can be formalized as modeling the probability over all possible items for user $u$ at time step $n_{u} {+} 1$:
\begin{equation*}
    p\bigl(v_{n_u+1}^{(u)}=v|\ \mathcal{S}_u\bigr)
\end{equation*}


\subsection{Model Architecture}

Here, we introduce a new sequential recommendation model called \textbf{BERT4Rec}, which adopts \textbf{B}idirectional \textbf{E}ncoder \textbf{R}epresentations from \textbf{T}ransformers to a new task, sequential \textbf{R}ecommendation.
It is built upon the popular self-attention layer, ``Transformer layer''. %

As illustrated in Figure~\ref{fig:model}b, BERT4Rec is stacked by $L$ bidirectional Transformer layers.
At each layer, it iteratively revises the representation of every position by exchanging information across all positions at the previous layer in parallel with the Transformer layer.
Instead of learning to pass relevant information forward step by step as RNN based methods did in Figure~\ref{fig:model}d, self-attention mechanism endows BERT4Rec with the capability to directly capture the dependencies in any distances.
This mechanism results in a global receptive field, while CNN based methods like Caser usually have a limited receptive field.
In addition, in contrast to RNN based methods, self-attention is straightforward to parallelize.




Comparing Figure~\ref{fig:model}b,~\ref{fig:model}c, and~\ref{fig:model}d, the most noticeable difference is that SASRec and RNN based methods are all left-to-right unidirectional architecture, while our BERT4Rec uses bidirectional self-attention to model users' behavior sequences.
In this way, our proposed model can obtain more powerful representations of users' behavior sequences to improve  recommendation performances.







\subsection{Transformer Layer}


As illustrated in Figure~\ref{fig:model}b, given an input sequence of length $t$, %
we iteratively compute hidden representations $\bm{h}_i^{l}$ at each layer $l$ for each position $i$ simultaneously by applying the Transformer layer from~\cite{Vaswani:nips2017:Attention}.
Here, we stack $\bm{h}_i^{l} \in \mathbb{R}^d$ together into matrix $\bm{H}^{l} {\in} \mathbb{R}^{t\times d}$ since we compute attention function on all positions simultaneously in practice. 
As shown in Figure~\ref{fig:model}a, the Transformer layer \texttt{Trm} contains two sub-layers, a \textit{Multi-Head Self-Attention} sub-layer and a \textit{Position-wise Feed-Forward Network}.

\textbf{Multi-Head Self-Attention}.
Attention mechanisms have become an integral part of sequence modeling in a variety of tasks, allowing capturing the dependencies between representation pairs without regard to their distance in the sequences.
Previous work has shown that it is beneficial to jointly attend to information from different representation subspaces at different positions~\cite{Vaswani:nips2017:Attention,Li:emnlp2018:Multi,Devlin:2018:BERT}.
Thus, we here adopt the multi-head self-attention instead of performing a single attention function.
Specifically, multi-head attention first linearly projects $\bm{H}^l$ into $h$ subspaces, with different, learnable linear projections, and then apply $h$ attention functions in parallel to produce the output representations which are concatenated and once again projected:
\begin{equation}
    \begin{aligned}
        \mathtt{MH}(\bm{H}^{l}) &= [\text{head}_1; \text{head}_2; \dots; \text{head}_h]\bm{W}^{O} \\
        \text{head}_i &= \mathtt{Attention}\bigl(\bm{H}^l\bm{W}_i^Q, \bm{H}^l\bm{W}_i^K, \bm{H}^l\bm{W}_i^V \bigr)
    \end{aligned}
\label{eq:mh}
\end{equation} %
where the projections matrices for each head $\bm{W}_i^Q \in \mathbb{R}^{d \times d/h}$, $\bm{W}_i^K \in \mathbb{R}^{d \times d/h}$, $\bm{W}_i^V \in \mathbb{R}^{d \times d/h}$, and $\bm{W}_i^O \in \mathbb{R}^{d \times d}$ are learnable parameters. %
Here, we omit the layer subscript $l$ for the sake of simplicity.
In fact, these projection parameters are not shared across the layers.
Here, the $\mathtt{Attention}$ function is \textit{Scaled Dot-Product Attention}: %
\begin{equation}
    \mathtt{Attention}(\bm{Q}, \bm{K}, \bm{V}) = \mathrm{softmax} \biggl(\frac{\bm{Q}\bm{K}^{\top}}{\sqrt{d/h}}\biggr) \bm{V}
    \label{eq:scale_dot}
\end{equation}
where \textit{query} $\bm{Q}$, \textit{key} $\bm{K}$, and \textit{value} $\bm{V}$ are projected from the same matrix $\bm{H}^l$ with different learned projection matrices as in Equation~\ref{eq:mh}. %
The \textit{temperature} $\sqrt{d/h}$ is introduced to produce a softer attention distribution for avoiding extremely small gradients~\cite{Hinton:nips2014:Distilling,Vaswani:nips2017:Attention}.


\textbf{Position-wise Feed-Forward Network}.
As described above, the self-attention sub-layer is mainly based on linear projections.
To endow the model with nonlinearity and interactions between different dimensions, we apply a \textit{Position-wise Feed-Forward Network} to the outputs of the self-attention sub-layer, separately and identically at each position.
It consists of two affine transformations with a Gaussian Error Linear Unit (\texttt{GELU}) activation in between:
\begin{equation}
    \begin{aligned}
      \texttt{PFFN}(\bm{H}^l) & = \bigl[\mathtt{FFN}(\bm{h}_1^{l})^{\top}; \dots; \mathtt{FFN}(\bm{h}_t^{l})^{\top}\bigr]^{\top}\\
       \mathtt{FFN}(\bm{x}) &= \texttt{GELU}\bigl(\bm{x} \bm{W}^{(1)}+\bm{b}^{(1)}\bigr)\bm{W}^{(2)} + \bm{b}^{(2)}\\
       \texttt{GELU}(x) &= x \Phi(x)
    \end{aligned}
    \label{eq:ffn}
\end{equation}
where $\Phi(x)$ is the cumulative distribution function of the standard gaussian distribution, $\bm{W}^{(1)}\in \mathbb{R}^{d\times 4d}$, $\bm{W}^{(2)}\in \mathbb{R}^{4d\times d}$, $\bm{b}^{(1)}\in \mathbb{R}^{4d}$ and $\bm{b}^{(2)}\in \mathbb{R}^{d}$ are learnable parameters and shared across all positions.
We omit the layer subscript $l$ for convenience.
In fact, these parameters are different from layer to layer.
In this work, following OpenAI GPT~\cite{Radford:2018:GPT} and BERT~\cite{Devlin:2018:BERT}, we use a smoother \texttt{GELU}~\cite{Hendrycks:gleu} activation rather than the standard \texttt{ReLu} activation. %


\textbf{Stacking Transformer Layer}.
As elaborated above, we can easily capture item-item interactions across the entire user behavior sequence using self-attention mechanism.
Nevertheless, it is usually beneficial to learn more complex item transition patterns by stacking the self-attention layers.
However, the network becomes more difficult to train as it goes deeper.
Therefore, we employ a residual connection~\cite{He:cvpr2016:resnet} around each of the two sub-layers as in Figure~\ref{fig:model}a, followed by layer normalization~\cite{Lei:Layer}.
Moreover, we also apply dropout~\cite{Srivastava:jmlr15:Dropout} to the output of each sub-layer, before it is normalized. %
That is, the output of each sub-layer is $\texttt{LN}(\bm{x} + \texttt{Dropout}(\text{sublayer}(\bm{x})))$, where $\text{sublayer}(\cdot)$ is the function implemented by the sub-layer itself, $\texttt{LN}$ is the layer normalization function defined in~\cite{Lei:Layer}. %
We use $\texttt{LN}$ to normalize the inputs over all the hidden units in the same layer for stabilizing and accelerating the network training.

In summary, BERT4Rec refines the hidden representations of each layer as follows:
\begin{align}
    \bm{H}^l &= \mathtt{Trm}\bigl(\bm{H}^{l-1}\bigr), \quad \forall i \in [1,\dots, L] \\
    \mathtt{Trm}(\bm{H}^{l-1}) &= \mathtt{LN}\Bigl(\bm{A}^{l-1} +  \mathtt{Dropout}\bigl(\mathtt{PFFN}(\bm{A}^{l-1})\bigr)\Bigr) \\
    \bm{A}^{l-1} &= \mathtt{LN}\Bigr(\bm{H}^{l-1} + \mathtt{Dropout}\bigl(\mathtt{MH}(\bm{H}^{l-1})\bigr)\Bigr) 
\end{align}


\subsection{Embedding Layer}
As elaborated above, without any recurrence or convolution module, the Transformer layer $\mathtt{Trm}$ is not aware of the order of the input sequence.
In order to make use of the sequential information of the input, we inject \textit{Positional Embeddings} into the input item embeddings at the bottoms of the Transformer layer stacks.
For a given item $v_i$, its input representation $\bm{h}_i^0$ is constructed by summing the corresponding item and positional embedding:
\begin{equation*}
    \begin{aligned}
        \bm{h}_{i}^0 &= \bm{v}_i + \bm{p}_i
    \end{aligned}
\end{equation*}
where $\bm{v}_i {\in} \bm{E}$ is the $d-$dimensional embedding for item $v_i$, $\bm{p}_i {\in} \bm{P}$ is the $d-$dimensional positional embedding for position index $i$.
In this work, we use the learnable positional embeddings instead of the fixed sinusoid embeddings in~\cite{Vaswani:nips2017:Attention} for better performances.
The positional embedding matrix $\bm{P} \in \mathbb{R}^{N\times d}$ allows our model to identify which portion of the input it is dealing with.
However, it also imposes a restriction on the maximum sentence length $N$ that our model can handle.
Thus, we need to truncate the the input sequence $[v_1,\dots,v_t]$ to the last $N$ items $[v_{t-N+1}^{u},\dots,v_t]$ if $t>N$.

\subsection{Output Layer}

After $L$ layers that hierarchically exchange information across all positions in the previous layer, we get the final output $\bm{H}^L$ for all items of the input sequence.
Assuming that we mask the item $v_t$ at time step $t$, we then predict the masked items $v_t$ base on $\bm{h}_t^L$ as shown in Figure~\ref{fig:model}b.
Specifically, we apply a two-layer feed-forward network with \texttt{GELU} activation in between to produce an output distribution over target items:
\begin{equation}
    P(v) = \mathrm{softmax}\bigl(\mathtt{GELU}(\bm{h}_t^L \bm{W}^P +\bm{b}^P) \bm{E}^{\top} + \bm{b}^O \bigr)
    \label{eq:prob}
\end{equation}
where $\bm{W}^P$ is the learnable projection matrix, $\bm{b}^P$, and $\bm{b}^O$ are bias terms, $\bm{E}\in \mathbb{R}^{|\mathcal{V}|\times d}$ is the embedding matrix for the item set $\mathcal{V}$.
We use the shared item embedding matrix in the input and output layer for alleviating overfitting and reducing model size.


\subsection{Model Learning}
\label{sec:ml}

\textbf{Training}. 
Conventional unidirectional sequential recommendation models usually train the model by predicting the next item for each position in the input sequence as illustrated in Figure~\ref{fig:model}c and~\ref{fig:model}d.
Specifically, the target of the input sequence $[v_1,\dots,v_{t}]$ is a shifted version $[v_2,\dots,v_{t+1}]$.
However, as shown in Figure~\ref{fig:model}b, jointly conditioning on both left and right context in a bidirectional model would cause the final output representation of each item to contain the information of the target item.
This makes predicting the future become trivial and the network would not learn anything useful.
A simple solution for this issue is to create $t-1$ samples (subsequences with next items like ([$v_1$], $v_2$) and ([$v_1,v_2$], $v_3$)) from the original length $t$ behavior sequence and then encode each historical subsequence with the bidirectional model to predict the target item. 
However, this approach is very time and resources consuming since we need to create a new sample for each position in the sequence and predict them separately.

In order to efficiently train our proposed model, we apply a new objective: \textit{Cloze} task~\cite{Wilson:Cloze} (also known as ``Masked Language Model'' in~\cite{Devlin:2018:BERT}) to sequential recommendation.
It is a test consisting of a portion of language with some words removed, where the participant is asked to fill the missing words.
In our case, for each training step, we randomly mask $\rho$ proportion of all items in the input sequence (\textit{i.e.}, replace with special token ``\texttt{[mask]}''), and then predict the original ids of the masked items based solely on its left and right context.
For example: 
\begin{figure}[h]
\centering
\resizebox{0.97\linewidth}{!}{
\begin{tikzpicture}[decoration={random steps,segment length=3mm, amplitude=0.5pt}]
\node[] (e1) at (0, 0) {\textbf{Input}: $[v_1,v_2,v_3,v_4,v_5]$};
\node[right of=e1, node distance=5.1cm] (e2)  {$[v_1,$ {\small$\mathtt{[mask]}_1$}$, v_3,$ {\small$\mathtt{[mask]}_2$}$, v_5]$};
\node[below of=e1, node distance=0.5cm, xshift=0.68cm] (e3)  {\textbf{Labels}: {\small$\mathtt{[mask]}_1$} $= v_2$, ~~ {\small$\mathtt{[mask]}_2$} $= v_4$};
\draw[FARROW] (e1) -> (e2) ;
\node[right of=e1, node distance=2.45cm, yshift=0.1cm] (e4)  {\scriptsize randomly mask};

\draw[decorate]  ($(e3.south west)+(-0.1, -0.1)$) rectangle ($(e2.north east)+(0.1, 0.1)$);

\end{tikzpicture}}
\end{figure}

\noindent The final hidden vectors corresponding to ``\texttt{[mask]}'' are fed into an output softmax over the item set, as in conventional sequential recommendation.
Eventually, we define the loss for each masked input $\mathcal{S}_u^{\prime}$ as the negative log-likelihood of the masked targets: %
\begin{equation}
    \mathcal{L} = \frac{1}{|\mathcal{S}_u^m|} \sum_{v_m\in \mathcal{S}_u^m} -\log P(v_m=v_m^{*}|\mathcal{S}_u^{\prime} )
    \label{eq:loss}
\end{equation}
where $\mathcal{S}_u^{\prime}$ is the masked version for user behavior history $\mathcal{S}_u$, $\mathcal{S}_u^m$ is the random masked items in it, $v_m^*$ is the true item for the masked item $v_m$, and the probability $P(\cdot)$ is defined in Equation~\ref{eq:prob}.

An additional advantage for Cloze task is that it can generate more samples to train the model.
Assuming a sequence of length $n$, conventional sequential predictions in Figure~\ref{fig:model}c and~\ref{fig:model}d  produce $n$ unique samples for training, while BERT4Rec can obtain $\binom{n}{k}$ samples (if we randomly mask $k$ items) in multiple epochs.
It allows us to train a more powerful bidirectional representation model.




\textbf{Test}.
As described above, we create a mismatch between the training and the final sequential recommendation task since the Cloze objective is to predict the current masked items while sequential recommendation aims to predict the future.
To address this, we append the special token ``\texttt{[mask]}'' to the end of user's behavior sequence, and then predict the next item based on the final hidden representation of this token.
To better match the sequential recommendation task (\textit{i.e.}, predict the last item), we also produce samples that only mask the last item in the input sequences during training.
It works like fine-tuning for sequential recommendation and can further improve the recommendation performances.

\subsection{Discussion}

Here, we discuss the relation of our model with previous related work.

\textbf{SASRec}. Obviously, SASRec is a left-to-right unidirectional version of our BERT4Rec with single head attention and causal attention mask.
Different architectures lead to different training methods.
SASRec predicts the next item for each position in a sequence, while BERT4Rec predicts the masked items in the sequence using Cloze objective.

\textbf{CBOW \& SG}. Another very similar work is Continuous Bag-of-Words (CBOW) and Skip-Gram (SG)~\cite{Mikolov:cbow}.
CBOW predicts a target word using the average of all the word vectors in its context (both left and right).
It can be seen as a simplified case of BERT4Rec, if we use one self-attention layer in BERT4Rec with uniform attention weights on items, unshare item embeddings, remove the positional embedding, and only mask the central item.
Similar to CBOW, SG can also be seen as a simplified case of BERT4Rec following similar reduction operations (mask all items except only one).
From this point of view, Cloze can be seen as a general form for the objective of CBOW and SG.
Besides, CBOW uses a simple aggregator to model word sequences since its goal is to learn good word representations, not sentence representations.
On the contrary, we seek to learn a powerful behavior sequence representation model (deep self-attention network in this work) for making recommendations.



\textbf{BERT}. Although our BERT4Rec is inspired by the BERT in NLP, it still has several differences from BERT:
\begin {enumerate*}[label=\bfseries\itshape\alph*\upshape)]
\item The most critical difference is that BERT4Rec is an end-to-end model for sequential recommendation, while BERT is a pre-training model for sentence representation.
BERT leverages large-scale task-independent corpora to pre-train the sentence representation model for various text sequence tasks since these tasks share the same background knowledge about the language.
However, this assumption does not hold in the recommendation tasks.
Thus we train BERT4Rec end-to-end for different sequential recommendation datasets.
\item Different from BERT, we remove the next sentence loss and segment embeddings since BERT4Rec models a user's historical behaviors as only one sequence in sequential recommendation task.
\end {enumerate*} 


\iffalse
\textbf{CBOW}. Another very similar work is Continuous Bag-of-Words (CBOW) model~\cite{Mikolov:cbow}.
CBOW predicts a target word using the average of all the word vectors in its context (both left and right).
It can be seen as a simplified case of BERT4Rec, if we replace the $L$-layer Transformer layer in BERT4Rec with uniform attention weights on items, use unshared item embeddings, remove the positional embedding, and only mask the central item.
CBOW uses a simple aggregator to model word sequences since its goal is to learn good word representations, not sentence representations.
On the contrary, we use a deep self-attention architecture to model the context since our goal is to learn a powerful sequence representation model for making recommendations.

\textbf{SG}. Similar to CBOW, Skip-Gram (SG)~\cite{Mikolov:nips2013:DRW} can also be seen as a simplified case of BERT4Rec following the similar reduction operations.
The difference is that we mask all items except only one for SG.
From this point of view, Cloze is a general form for the objective of CBOW and SG.

\textbf{FISM}. If we use one self-attention layer with uniform attention weights on items, use unshared item embeddings, and
 remove the position embedding, SASRec is reduced to FISM. 
Thus our model can be viewed as an adaptive, hierarchical, sequential item similarity model for next item recommendation.
\fi






